\begin{abstract}
Multi-agent applications accumulate reusable context while executing
interdependent LLM requests. A worker failure can lose this state and delay
several branches at once, yet protecting every context in fast memory is
expensive. We propose AegisKV, a workflow-aware manager for KV checkpoints
across Peer HBM, shared DRAM and shared storage. It retains complete
checkpoints in DRAM or storage and replicates selected layers and hot
blocks on a preassigned takeover DP. AegisKV values Peer acceleration,
whole-checkpoint demotion and final eviction separately, combining normal reuse
with the workflow cost of recovery. It evaluates requests affected by a
common failure together, uses sparse demand and layer overlap to value
fast-tier residency, and directly re-evaluates complete replacement plans
under bounded decision work and a migration threshold.
Versioned transitions preserve recovery validity during migration and
explicit fallback. A per-DP Adapter limits Peer capacity while the Master
chooses per-request layer depth and hot blocks without changing native
attention semantics. We connect trace-derived
workflow observations and existing reconstruction microbenchmarks to a
reproducible analytical study and a complete experimental protocol.
In 36 declared Peer-disabled scenarios, the placement model improves on
LRU in 17 and ties in 19. Matched analytical schedules illustrate Peer's
incremental overlap benefit; neither model measures a three-tier system. Historical outage timing drives a separate recovery-load model in which
a lower mean coexists with worse tail latency. Mechanism ablations expose
interactions between layer placement and overlap, and a rejected rental
alternative. The study identifies when contention, prediction error, and commit
lag limit the benefit. End-to-end fault injection and normal
service measurements remain necessary to establish the proposed system's
performance.
\end{abstract}
